% =====================================================================
%  سورة أم القرآن — الأبيات ١٠٨–١١٥
% =====================================================================
\chapter{سورة أم القرآن}

يذكر الناظم في هذا الباب خلاف القراء في سورة الفاتحة (وتسمى أم القرآن)، ثم يجعل بعض كلماتها أصولًا مطّردة في القرآن كله: كلمة \textbf{الصراط}، وهاء \textbf{عليهم وإليهم ولديهم}، و\textbf{ميم الجمع}.
وعدد أبياته ثمانية.

\begin{tanbih}[قبل أن تبدأ: كيف تقرأ رموز هذا الباب]
الرمز يأتي \textbf{بعد} الكلمة القرآنية. فانظر إلى أول حرف من الكلمة التي تأتي بعدها: إن كان من حروف «أبجد هوز...» فهو رمز قارئ (راجع جدول الرموز في دليل الاستخدام).
وفي الأبيات ظلّلنا الكلمة القرآنية المختلف فيها بالأصفر \qw{هكذا}، ولوّنّا كلمة الرمز بلون القارئ.
\end{tanbih}

% ---------------------------------------------------------------------
\section{﴿مالك يوم الدين﴾ و﴿الصراط﴾}
% ---------------------------------------------------------------------
\bayt{108}{\qw{وَمَالِكِ} يَوْمِ الدِّينِ \rz{Kisai}{رَاوِيهِ} \rz{Asim}{نَاصِرٌ}}{وَعِنْدَ \qw{سِرَاطٍ} \qw{وَالسِّرَاطَ} \qname{IbnKathir}{لِ قُنْبُلَا}}

\begin{mufradat}
\kalima{رَاوِيهِ}{الراء رمز \Kisai. ومعنى الكلمة: ناقله.}
\kalima{نَاصِرٌ}{النون رمز \Asim. ومعناها: مؤيِّده.}
\kalima{لِ قُنْبُلَا}{«لِ» فعل أمر من «وَلِيَ يَلِي» أي: اتْبَعْ قنبلًا. وليست اللام هنا رمز \Hisham{}؛ لأنه صرّح بعدها باسم \Qunbul.}
\end{mufradat}

\maana
قرأ \Kisai{} و\Asim{} \ayah{مَالِكِ يَوْمِ الدِّينِ} بإثبات الألف بعد الميم، ولم يحتج الناظم إلى أن يقول «بالمد» لأن لفظه بالكلمة «مالك» يدل على القراءة (وهذا معنى قوله في المقدمة: «وباللفظ أستغني عن القيد إن جلا»).
فتكون قراءة الباقين بحذف الألف: \ayah{مَلِكِ يَوْمِ الدِّينِ}.

ثم قال: اتبع \Qunbul{} في قراءة \textbf{سراط} و\textbf{السراط} بالسين حيث وقع في القرآن.

\begin{qiraat}
\qira{\Asim{} و\Kisai}{\ayah{مَالِكِ} بالألف}
\qira{\Baqun{} (\Nafi، \IbnKathir، \AbuAmr، \IbnAmir، \Hamza)}{\ayah{مَلِكِ} بلا ألف}
\qira{\Qunbul}{\ayah{السِّرَاطَ} و\ayah{سِرَاطَ} بالسين في كل القرآن}
\end{qiraat}

\begin{qaida}[قاعدة كلمة «صراط» عند قنبل]
تقرأ لقنبل بالسين في كل موضع، سواء كانت الكلمة:
\begin{itemize}
\item \textbf{معرّفة بأل}: \ayah{ٱهۡدِنَا ٱلسِّرَٰطَ ٱلۡمُسۡتَقِيمَ} \src{الفاتحة: ٦}
\item \textbf{نكرة}: \ayah{وَإِنَّكَ لَتَهۡدِيٓ إِلَىٰ سِرَٰطࣲ مُّسۡتَقِيمࣲ} \src{الشورى: ٥٢}
\item \textbf{مضافة}: \ayah{سِرَٰطَ ٱلَّذِينَ أَنۡعَمۡتَ عَلَيۡهِمۡ} \src{الفاتحة: ٧}، \ayah{وَأَنَّ هَٰذَا سِرَٰطِي مُسۡتَقِيمࣰا} \src{الأنعام: ١٥٣}
\end{itemize}
وكلمة «حيث أتى» في أول البيت التالي هي التي أفادت هذا التعميم.
\end{qaida}

% ---------------------------------------------------------------------
\bayt{109}{بِحَيْثُ أَتَى وَالصَّادَ زَايًا أَشِمَّهَا}{لَدَى \qname{Hamza}{خَلَفٍ} وَاشْمِمْ \qname{Hamza}{لِخَلَّادٍ} الْاَوَّلَا}

\begin{mufradat}
\kalima{بحيث أتى}{أي في كل موضع جاءت فيه الكلمة في القرآن.}
\kalima{أَشِمَّها زايًا}{الإشمام هنا: أن تخلط صوت الصاد بصوت الزاي.}
\kalima{الأوّلا}{الموضع الأول فقط، وهو \ayah{ٱهۡدِنَا ٱلصِّرَٰطَ ٱلۡمُسۡتَقِيمَ} في الفاتحة.}
\end{mufradat}

\maana
قرأ \Khalaf{} بإشمام الصاد صوت الزاي في «صراط» و«الصراط» حيث وقع، وقرأ \Khallad{} بالإشمام في الموضع الأول فقط وهو \ayah{ٱهۡدِنَا ٱلصِّرَٰطَ ٱلۡمُسۡتَقِيمَ}، وما سواه بالصاد الخالصة.
وقرأ الباقون (ومعهم \Khallad{} في غير الموضع الأول) بالصاد الخالصة.

\begin{qiraat}
\qira{\Qunbul}{بالسين: \ayah{ٱلسِّرَٰطَ}}
\qira{\Khalaf}{بالإشمام (صاد مخلوطة بالزاي) في كل القرآن}
\qira{\Khallad}{بالإشمام في \ayah{ٱهۡدِنَا ٱلصِّرَٰطَ} فقط، والباقي بالصاد}
\qira{\Baqun}{بالصاد الخالصة}
\end{qiraat}

\begin{qaida}[كيف يُنطق الإشمام هنا؟]
تخلط لفظ الصاد بلفظ الزاي فيتولد حرف ليس صادًا خالصًا ولا زايًا خالصًا، لكن صوت الصاد هو الغالب.
وتقريبه: أن تنطق الصاد قريبًا من نطق العامّة للظاء (كما قال صاحب الوافي).
\end{qaida}

\begin{tanbih}[تنبيه من الوافي]
ذكر الشيخ القاضي في حاشية الوافي أن \Khallad{} في الموضع الأول (\ayah{ٱهۡدِنَا ٱلصِّرَٰطَ}) له وجهان: الإشمام والصاد الخالصة، وأن اقتصار الناظم على الإشمام فيه قصور. فاعرض ذلك على شيخك عند الأداء.
\end{tanbih}

% ---------------------------------------------------------------------
\section{هاء ﴿عليهم﴾ و﴿إليهم﴾ و﴿لديهم﴾}
% ---------------------------------------------------------------------
\bayt{110}{\qw{عَلَيْهِمْ} \qw{إِلَيْهِمْ} \qname{Hamza}{حَمْزَةٌ} \qw{وَلَدَيْهِمُو}}{جَمِيعًا بِضَمِّ الْهَاءِ وَقْفًا وَمَوْصِلَا}

\begin{mufradat}
\kalima{وقفًا وموصلًا}{أي في حال الوقف وفي حال الوصل.}
\end{mufradat}

\maana
قرأ \Hamza{} هذه الكلمات الثلاث فقط: \textbf{عليهم، إليهم، لديهم} بضم الهاء في كل القرآن، وقفًا ووصلًا، سواء جاء بعدها متحرك أو ساكن.
وقرأ الباقون بكسر الهاء (وأُخذ الكسر من لفظ الناظم بالكلمات مكسورة).

\begin{qiraat}
\qira{\Hamza}{\ayah{عَلَيۡهُمۡ}، \ayah{إِلَيۡهُمۡ}، \ayah{لَدَيۡهُمۡ} بضم الهاء}
\qira{\Baqun}{\ayah{عَلَيۡهِمۡ}، \ayah{إِلَيۡهِمۡ}، \ayah{لَدَيۡهِمۡ} بكسر الهاء}
\end{qiraat}

\begin{amthila}
\begin{itemize}
\item بعدها متحرك: \ayah{غَيۡرِ ٱلۡمَغۡضُوبِ عَلَيۡهِمۡ وَلَا ٱلضَّآلِّينَ} \src{الفاتحة: ٧} \ila{} \Hamza: \textbf{عليهُمْ}.
\item \ayah{يُحِبُّونَ مَنۡ هَاجَرَ إِلَيۡهِمۡ} \src{الحشر: ٩}، \ayah{وَمَا كُنتَ لَدَيۡهِمۡ إِذۡ يُلۡقُونَ} \src{آل عمران: ٤٤}.
\item بعدها ساكن: \ayah{حَقَّ عَلَيۡهِمُ ٱلۡقَوۡلُ} \src{القصص: ٦٣}، \ayah{إِذۡ أَرۡسَلۡنَآ إِلَيۡهِمُ ٱثۡنَيۡنِ} \src{يس: ١٤} \ila{} \Hamza{} يضم الهاء هنا أيضًا.
\end{itemize}
\end{amthila}

\begin{tanbih}
الحكم خاص بهذه الكلمات الثلاث بلفظها. أما نحو \textbf{فيهم، بهم، يزكيهم} فلا يضم \Hamza{} هاءها، بل يكسرها كالباقين (وسيأتي حكمها عند ساكن في البيت ١١٤).
\end{tanbih}

% ---------------------------------------------------------------------
\section{ميم الجمع}
% ---------------------------------------------------------------------
\begin{qaida}[ما هي ميم الجمع؟ وما معنى صلتها؟]
\textbf{ميم الجمع}: الميم الزائدة الدالة على جمع الذكور، مثل: عليه\textbf{م}، أنت\textbf{م}، ك\textbf{م}.\\
\textbf{الصلة}: أن تضم الميم وتمدها بواو لفظية مقدار حركتين وصلًا: \textbf{عليهمُو}، \textbf{أنتمُو}.\\
\textbf{الإسكان}: أن تنطقها ساكنة: \textbf{عليهمْ}.\\
والخلاف كله \textbf{في الوصل فقط}. أما عند الوقف عليها فالجميع يسكّنها.
\end{qaida}

\bayt{111}{وَصِلْ ضَمَّ مِيمِ الْجَمْعِ قَبْلَ مُحَرَّكٍ}{\rz{IbnKathir}{دِرَاكًا} \qname{Nafi}{وَقَالُونٌ} بِتَخْيِيرِهِ جَلَا}

\begin{mufradat}
\kalima{قبل محرك}{أي إذا جاء بعد الميم حرف متحرك.}
\kalima{دِرَاكًا}{متتابعًا. والدال رمز \IbnKathir.}
\kalima{جَلَا}{كشف وبيّن. والجيم هنا \textbf{ليست} رمزًا لورش؛ لأنه صرّح باسم \Qalun{} قبلها.}
\end{mufradat}

\maana
إذا وقعت ميم الجمع قبل حرف متحرك فـ\IbnKathir{} يضمها ويصلها بواو في كل القرآن، و\Qalun{} مخيَّر بين الصلة والإسكان (فله وجهان).

\begin{amthila}
\ayah{أَنۡعَمۡتَ عَلَيۡهِمۡ غَيۡرِ ٱلۡمَغۡضُوبِ} \src{الفاتحة: ٧} \ila{} \IbnKathir{} (ووجه لـ\Qalun): \textbf{عليهِمُو غيرِ}.\\
\ayah{ءَأَنذَرۡتَهُمۡ أَمۡ لَمۡ تُنذِرۡهُمۡ} \src{البقرة: ٦} \ila{} \textbf{ءأنذرتهمُو أم لم تنذرهمُو لا يؤمنون}.
\end{amthila}

% ---------------------------------------------------------------------
\bayt{112}{وَمِنْ قَبْلِ هَمْزِ الْقَطْعِ صِلْهَا \qname{Nafi}{لِوَرْشِهِمْ}}{وَأَسْكَنَهَا الْبَاقُونَ بَعْدُ لِتَكْمُلَا}

\maana
\Warsh{} يصل ميم الجمع بواو \textbf{إذا كان بعدها همزة قطع} فقط (وتُمدّ عنده مدًّا منفصلًا مشبعًا). وأسكنها باقي القراء بعد \IbnKathir{} و\Qalun{} و\Warsh{}، أي: \AbuAmr{} و\IbnAmir{} و\Asim{} و\Hamza{} و\Kisai.
ونصّ الناظم على الإسكان لأن الإسكان لا يؤخذ من الضد.

\begin{qiraat}
\qira{\IbnKathir}{الصلة دائمًا قبل أي متحرك}
\qira{\Qalun}{وجهان: الصلة والإسكان قبل أي متحرك}
\qira{\Warsh}{الصلة قبل همزة القطع فقط، والإسكان قبل غيرها}
\qira{\AbuAmr، \IbnAmir، \Asim، \Hamza، \Kisai}{الإسكان دائمًا}
\end{qiraat}

\begin{amthila}
\ayah{عَلَيۡكُمۡ أَنفُسَكُمۡ} \src{المائدة: ١٠٥}، \ayah{وَمِنۡهُمۡ أُمِّيُّونَ} \src{البقرة: ٧٨} \ila{} \Warsh: \textbf{عليكمُو أنفسكم}، \textbf{ومنهمُو أميون} (مع المد المشبع).\\
أما \ayah{عَلَيۡهِمۡ غَيۡرِ} فـ\Warsh{} يسكّنها لأن بعدها غير همزة.
\end{amthila}

% ---------------------------------------------------------------------
\bayt{113}{وَمِنْ دُونِ وَصْلٍ ضَمُّهَا قَبْلَ سَاكِنٍ}{لِكُلٍّ وَبَعْدَ الْهَاءِ كَسْرُ \rz{AbuAmr}{فَتَى الْعَلَا}}

\begin{mufradat}
\kalima{فتى العلا}{\AbuAmr{} البصري، وأبوه اسمه العلاء.}
\end{mufradat}

\maana
إذا جاء بعد ميم الجمع \textbf{ساكن} فكل القراء يضمونها \textbf{بلا صلة} (لأن الصلة تُحذف لالتقاء الساكنين)، نحو: \ayah{عَلَيۡكُمُ ٱلصِّيَامُ} \src{البقرة: ١٨٣}، \ayah{وَمِنۡهُمُ ٱلَّذِينَ} \src{التوبة: ٦١}.

ثم استثنى \AbuAmr{}: فهو يكسر الميم إذا كانت بعد هاء، بشرط مذكور في البيت التالي.

\bayt{114}{مَعَ الْكَسْرِ قَبْلَ الْهَا أَوِ الْيَاءِ سَاكِنًا}{وَفِي الْوَصْلِ كَسْرُ الْهَاءِ بِالضَّمِّ \rzg{شَمْلَلَا}}

\begin{mufradat}
\kalima{شَمْلَلَا}{أسرع. والشين رمز \Hamza{} و\Kisai{} معًا.}
\end{mufradat}

\bayt{115}{كَمَا بِهِمُ الْأَسْبَابُ ثُمَّ عَلَيْهِمُ الْـ}{ـقِتَالُ وَقِفْ لِلْكُلِّ بِالْكَسْرِ مُكْمِلَا}

\begin{mufradat}
\kalima{كما}{الكاف هنا ليست رمزًا لابن عامر، بل هي كاف التشبيه وبعدها مثال (وهذا من الحروف التي «لا ريبة في اتصالها» كما قال في المقدمة).}
\kalima{مُكْمِلًا}{متمِّمًا للحكم.}
\end{mufradat}

\maana[(الأبيات ١١٣–١١٥)]
هذه حالة خاصة: ميم جمع \textbf{بعد هاء}، وقبل الهاء \textbf{كسرة أو ياء ساكنة}، وبعد الميم \textbf{ساكن}. ففيها ثلاث قراءات في الوصل:

\begin{qiraat}
\qira{\AbuAmr}{كسر الهاء وكسر الميم: \textbf{بِهِمِ الأسبابُ}، \textbf{عليهِمِ القتالُ}}
\qira{\Hamza{} و\Kisai}{ضم الهاء وضم الميم: \textbf{بِهُمُ الأسبابُ}، \textbf{عليهُمُ القتالُ}}
\qira{\Baqun}{كسر الهاء وضم الميم: \textbf{بِهِمُ الأسبابُ}، \textbf{عليهِمُ القتالُ}}
\end{qiraat}

\begin{amthila}
\begin{itemize}
\item قبل الهاء كسرة: \ayah{وَتَقَطَّعَتۡ بِهِمُ ٱلۡأَسۡبَابُ} \src{البقرة: ١٦٦}، \ayah{فِي قُلُوبِهِمُ ٱلۡعِجۡلَ} \src{البقرة: ٩٣}.
\item قبل الهاء ياء ساكنة: \ayah{كُتِبَ عَلَيۡهِمُ ٱلۡقِتَالُ} \src{البقرة: ٢٤٦}، \ayah{يُرِيهِمُ ٱللَّهُ أَعۡمَٰلَهُمۡ} \src{البقرة: ١٦٧}، \ayah{إِلَيۡهِمُ ٱثۡنَيۡنِ} \src{يس: ١٤}.
\end{itemize}
\end{amthila}

\begin{qaida}[حكم الوقف على هذه الكلمات]
عند الوقف على الميم تُسكَّن لكل القراء، وتُكسر الهاء لكل القراء (لأن سبب ضمها أو كسرها عند بعضهم هو الساكن الذي بعدها، وقد زال بالوقف): \textbf{بِهِمْ}، \textbf{يُرِيهِمْ}.
\textbf{ويُستثنى} \Hamza{} في الكلمات الثلاث \textbf{عليهم، إليهم، لديهم}؛ فإنه يضم هاءها وقفًا ووصلًا كما سبق في البيت ١١٠: \textbf{عليهُمْ}.
\end{qaida}

\begin{khulasa}[جدول جامع: ﴿عليهم القتال﴾ و﴿يريهم الله﴾]
\centering
\renewcommand{\arraystretch}{1.5}
\begin{tabular}{l c c c}
\toprule
\textbf{القارئ} & \textbf{عليهم القتال (وصلًا)} & \textbf{يريهم الله (وصلًا)} & \textbf{الوقف} \\
\midrule
\AbuAmr & عليهِمِ & يريهِمِ & عليهِمْ / يريهِمْ \\
\Hamza & عليهُمُ & يريهُمُ & عليهُمْ / يريهِمْ \\
\Kisai & عليهُمُ & يريهُمُ & عليهِمْ / يريهِمْ \\
\Baqun & عليهِمُ & يريهِمُ & عليهِمْ / يريهِمْ \\
\bottomrule
\end{tabular}
\end{khulasa}

% ---------------------------------------------------------------------
\section{خلاصة الباب}
% ---------------------------------------------------------------------
\begin{khulasa}[خلاصة سورة أم القرآن]
\renewcommand{\arraystretch}{1.45}
\begin{tabularx}{\linewidth}{>{\raggedleft\arraybackslash}p{0.2\linewidth} X}
\toprule
\textbf{الكلمة / الأصل} & \textbf{القراءات} \\
\midrule
\ayah{مَٰلِكِ} & \Asim{} و\Kisai: بالألف. \Baqun: \textbf{مَلِكِ}. \\
\ayah{ٱلصِّرَٰطَ} & \Qunbul: بالسين. \Khalaf: بالإشمام كل القرآن. \Khallad: الإشمام في الفاتحة فقط. \Baqun: بالصاد. \\
\ayah{عَلَيۡهِمۡ} وأختاها & \Hamza: بضم الهاء وقفًا ووصلًا. \Baqun: بالكسر. \\
ميم الجمع قبل متحرك & \IbnKathir: صلة. \Qalun: صلة وإسكان. \Warsh: صلة قبل الهمز فقط. \Baqun: إسكان. \\
ميم الجمع قبل ساكن & الكل بالضم بلا صلة، إلا ما بعد هاء قبلها كسر أو ياء: \AbuAmr{} يكسر الميم، و\Hamza{} و\Kisai{} يضمان الهاء والميم. \\
\bottomrule
\end{tabularx}
\end{khulasa}
